\documentclass[10pt, letterpaper]{article}

% Packages:
\usepackage[
    ignoreheadfoot, % set margins without considering header and footer
    top=2 cm, % seperation between body and page edge from the top
    bottom=2 cm, % seperation between body and page edge from the bottom
    left=2 cm, % seperation between body and page edge from the left
    right=2 cm, % seperation between body and page edge from the right
    footskip=1.0 cm, % seperation between body and footer
    % showframe % for debugging 
]{geometry} % for adjusting page geometry
\usepackage{titlesec} % for customizing section titles
\usepackage{tabularx} % for making tables with fixed width columns
\usepackage{array} % tabularx requires this
\usepackage[dvipsnames]{xcolor} % for coloring text
\definecolor{primaryColor}{RGB}{0, 79, 144} % define primary color
\usepackage{enumitem} % for customizing lists
\usepackage{fontawesome5} % for using icons
\usepackage{amsmath} % for math
\usepackage[
    pdftitle={Xu Han's CV},
    pdfauthor={Xu Han},
    pdfcreator={LaTeX with RenderCV},
    colorlinks=true,
    urlcolor=primaryColor
]{hyperref} % for links, metadata and bookmarks
\usepackage[pscoord]{eso-pic} % for floating text on the page
\usepackage{calc} % for calculating lengths
\usepackage{bookmark} % for bookmarks
\usepackage{lastpage} % for getting the total number of pages
\usepackage{changepage} % for one column entries (adjustwidth environment)
\usepackage{paracol} % for two and three column entries
\usepackage{ifthen} % for conditional statements
\usepackage{needspace} % for avoiding page brake right after the section title
\usepackage{etoolbox}
\usepackage{iftex} % check if engine is pdflatex, xetex or luatex

% Ensure that generate pdf is machine readable/ATS parsable:
\ifPDFTeX
    \input{glyphtounicode}
    \pdfgentounicode=1
    % \usepackage[T1]{fontenc} % this breaks sb2nov
    \usepackage[utf8]{inputenc}
    \usepackage{lmodern}
\fi



% Some settings:
\AtBeginEnvironment{adjustwidth}{\partopsep0pt} % remove space before adjustwidth environment
\pagestyle{empty} % no header or footer
\setcounter{secnumdepth}{0} % no section numbering
\setlength{\parindent}{0pt} % no indentation
\setlength{\topskip}{0pt} % no top skip
\setlength{\columnsep}{0cm} % set column seperation
\makeatletter
\let\ps@customFooterStyle\ps@plain % Copy the plain style to customFooterStyle
\patchcmd{\ps@customFooterStyle}{\thepage}{
    \color{gray}\textit{\small Xu Han - Page \thepage{} of \pageref*{LastPage}}
}{}{} % replace number by desired string
\makeatother
\pagestyle{customFooterStyle}

\titleformat{\section}{\needspace{4\baselineskip}\bfseries\large}{}{0pt}{}[\vspace{1pt}\titlerule]

\titlespacing{\section}{
    % left space:
    -1pt
}{
    % top space:
    0.3 cm
}{
    % bottom space:
    0.2 cm
} % section title spacing

\renewcommand\labelitemi{$\circ$} % custom bullet points
\newenvironment{highlights}{
    \begin{itemize}[
        topsep=0.10 cm,
        parsep=0.10 cm,
        partopsep=0pt,
        itemsep=0pt,
        leftmargin=0.4 cm + 10pt
    ]
}{
    \end{itemize}
} % new environment for highlights

\newenvironment{highlightsforbulletentries}{
    \begin{itemize}[
        topsep=0.10 cm,
        parsep=0.10 cm,
        partopsep=0pt,
        itemsep=0pt,
        leftmargin=10pt
    ]
}{
    \end{itemize}
} % new environment for highlights for bullet entries


\newenvironment{onecolentry}{
    \begin{adjustwidth}{
        0.2 cm + 0.00001 cm
    }{
        0.2 cm + 0.00001 cm
    }
}{
    \end{adjustwidth}
} % new environment for one column entries

\newenvironment{twocolentry}[2][]{
    \onecolentry
    \def\secondColumn{#2}
    \setcolumnwidth{\fill, 4.5 cm}
    \begin{paracol}{2}
}{
    \switchcolumn \raggedleft \secondColumn
    \end{paracol}
    \endonecolentry
} % new environment for two column entries

\newenvironment{header}{
    \setlength{\topsep}{0pt}\par\kern\topsep\centering\linespread{1.5}
}{
    \par\kern\topsep
} % new environment for the header

\newcommand{\placelastupdatedtext}{% \placetextbox{<horizontal pos>}{<vertical pos>}{<stuff>}
  \AddToShipoutPictureFG*{% Add <stuff> to current page foreground
    \put(
        \LenToUnit{\paperwidth-2 cm-0.2 cm+0.05cm},
        \LenToUnit{\paperheight-1.0 cm}
    ){\vtop{{\null}\makebox[0pt][c]{
        \small\color{gray}\textit{Updated 27 Oct 2026}\hspace{\widthof{Updated 27 Oct 2026}}
    }}}%
  }%
}%

% save the original href command in a new command:
\let\hrefWithoutArrow\href

% new command for external links:
\renewcommand{\href}[2]{\hrefWithoutArrow{#1}{\ifthenelse{\equal{#2}{}}{ }{#2\nobreakspace}\raisebox{.15ex}{\footnotesize \faExternalLink*}}}


\begin{document}
\newcommand{\AND}{\unskip\cleaders\copy\ANDbox\hskip\wd\ANDbox\ignorespaces}
\newsavebox\ANDbox
\sbox\ANDbox{}
\placelastupdatedtext
\begin{header}
\textbf{\fontsize{24 pt}{24 pt}\selectfont Xu Han}
\par
\vspace{0.3 cm}
\normalsize
\mbox{{\color{black}\footnotesize\faMapMarker*}\hspace*{0.13cm}Munich, Germany}\kern 0.25 cm
\mbox{\hrefWithoutArrow{mailto:xuhanphysik@gmail.com}{\color{black}{\footnotesize\faEnvelope[regular]}\hspace*{0.13cm}xuhanphysik@gmail.com}}\kern 0.25 cm
\mbox{\hrefWithoutArrow{https://github.com/xhvoid}{\color{black}{\footnotesize\faGithub}\hspace*{0.13cm}xhvoid}}
\end{header}
\vspace{0.10 cm}
\begin{onecolentry}
\centering\hrefWithoutArrow{https://xhvoid.github.io/Xu-Han.github.io/}{\footnotesize\faLink\hspace*{0.13cm}xhvoid.github.io/Xu-Han.github.io}
\end{onecolentry}

\vspace{0.10 cm}
\begin{onecolentry}
Photonics M.Sc. student and LMU astrophysics graduate focused on astronomical instrumentation, wavefront sensing and optical-system performance, combining optical modelling, experimental apparatus development and astronomical data analysis.
\end{onecolentry}
\vspace{0.05 cm}
\section{Education}
\begin{twocolentry}{\textit{Mar 2026 - Present}}
\textbf{M.Sc. Photonics} | Munich University of Applied Sciences
\end{twocolentry}
\begin{onecolentry}
\textbf{Selected coursework:} Physical Optics and Modelling (\textbf{Zemax optical design}); Micro and Fiber Optics; Optical Measurement Technology; Semiconductor Optics and Components.
\end{onecolentry}
\vspace{0.08 cm}
\begin{twocolentry}{\textit{Apr 2022 - Aug 2024}}
\textbf{M.Sc. Astrophysics} | LMU Munich
\end{twocolentry}
\begin{onecolentry}
Grade: 1.55 (German scale). \textbf{Selected coursework:} Essentials of Advanced Astrophysics; Cosmology and Structure Formation; Statistical Methods; Instrumental Laboratory Courses.
\end{onecolentry}
\vspace{0.08 cm}
\begin{twocolentry}{\textit{Oct 2018 - Mar 2022}}
\textbf{B.Sc. Physics} | University of Duisburg-Essen
\end{twocolentry}
\begin{onecolentry}
Grade: 1.8 (German scale).
\end{onecolentry}
\vspace{0.08 cm}
\begin{twocolentry}{\textit{Sep 2013 - Jul 2017}}
\textbf{B.Eng. Electrical Engineering and Automation}

Hubei University of Technology
\end{twocolentry}
\vspace{0.08 cm}
\section{Selected Research and Engineering Projects}
\begin{onecolentry}
{\large\textbf{Shack-Hartmann Wavefront Sensing \& Adaptive Optics Simulation}}
\end{onecolentry}
\begin{twocolentry}{\textit{Mar 2026 - Present}}
\textit{Independent project (ongoing development)}
\end{twocolentry}
\begin{onecolentry}
\hrefWithoutArrow{https://xhvoid.github.io/Shack-Hartmann-AO-Simulation/}{Interactive explorer} | \hrefWithoutArrow{https://github.com/xhvoid/Shack-Hartmann-AO-Simulation}{Source code}
\end{onecolentry}
\begin{onecolentry}
\begin{highlights}
\item Developed a modular Python SCAO framework connecting atmospheric turbulence, detector-level Shack-Hartmann sensing, calibrated reconstruction, deformable-mirror correction and delayed feedback.
\item Investigated detector-noise and gain/latency trade-offs using residual wavefront error, propagated-PSF Strehl and encircled energy.
\item Diagnosed a reconstruction failure associated with measurement dropout in a 2.2 m NGS model; stronger TSVD regularisation kept reconstruction usable throughout the matched two-second run.
\item Cross-checked selected optical components against HCIPy and analytic references; documented reproducible experiments, model limitations and an interactive explorer of recorded results.
\end{highlights}
\end{onecolentry}
\vspace{0.08 cm}
\begin{twocolentry}{\textit{Mar 2021 - Aug 2021}}
{\large\textbf{\href{https://xhvoid.github.io/Xu-Han.github.io/projects/\#experimental-apparatus}{Experimental Collisional Charging Apparatus}}}

\textit{Bachelor's thesis | University of Duisburg-Essen}
\end{twocolentry}
\begin{onecolentry}
Supervisors: Prof. Gerhard Wurm and Dr. Felix Jungmann.
\end{onecolentry}
\begin{onecolentry}
\begin{highlights}
\item Designed the mechanical apparatus in CAD, fabricated components by 3D printing and assembled the setup; integrated and debugged electronic charge-measurement modules.
\item Conducted measurements and analysed grain-size dependence of collisional charging, testing a dust-collision growth hypothesis relevant to early planet formation.
\end{highlights}
\end{onecolentry}
\vspace{0.08 cm}
\begin{twocolentry}{\textit{Mar 2023 - Aug 2024}}
{\large\textbf{\href{https://xhvoid.github.io/Xu-Han.github.io/files/Watershed\_Voids\%EF\%BC\%88thesis\%EF\%BC\%89.pdf}{Weak Lensing Void Finder}}}

\textit{Master's thesis | LMU Munich}
\end{twocolentry}
\begin{onecolentry}
Supervisors: Prof. Joseph Mohr and Dr. Chris Davies.
\end{onecolentry}
\begin{onecolentry}
\begin{highlights}
\item Developed and parameter-tested a watershed finder for cosmic voids in 2D full-sky weak-lensing convergence maps.
\item Applied the analysis to DES Y3 weak-lensing maps; characterised void abundance, effective radii, ellipticities and convergence profiles.
\end{highlights}
\end{onecolentry}
\vspace{0.08 cm}
\clearpage
\section{Laboratory and Instrumentation Training}
\begin{onecolentry}
\begin{highlights}
\item \textbf{Photonics laboratory:} Completed experiments: APD measurements, laser spectra, nonlinear optics, laser interferometry, fiber Bragg gratings and fiber coupling.
\item \textbf{Astronomical instrumentation laboratory:} Completed experiments: CCD imaging and data reduction, optical interferometry, deformable-mirror AO tuning and radio measurements of Galactic rotation.
\item \textbf{Observatory experience:} Participated in observing and image processing at Wendelstein Observatory using the 43 cm telescope.
\end{highlights}
\end{onecolentry}
\section{Additional Independent Projects}
\begin{onecolentry}
{\large\textbf{\href{https://github.com/xhvoid/Astrophotonics}{Astrophotonics}}}
\end{onecolentry}
\begin{onecolentry}
\begin{highlights}
\item Modelled how AO wavefront error, injection jitter and photonic-lantern mode count affect fiber coupling, spectral SNR and radial-velocity precision.
\end{highlights}
\end{onecolentry}
\vspace{0.08 cm}
\begin{onecolentry}
{\large\textbf{\href{https://github.com/xhvoid/AO-Fed-IFU-Spectrograph-Simulation}{AO-Fed IFU Spectrograph Simulation}}}
\end{onecolentry}
\begin{onecolentry}
\begin{highlights}
\item Modelled the optical-to-detector chain and datacube reconstruction; compared recovered flux, velocity and SNR with synthetic truth and examined calibration sensitivity.
\end{highlights}
\end{onecolentry}
\vspace{0.08 cm}
\begin{onecolentry}
{\large\textbf{\href{https://github.com/xhvoid/Optoelectronic-Signal-Chain-Simulations}{Optoelectronic Signal Chain Simulations}}}
\end{onecolentry}
\begin{onecolentry}
\begin{highlights}
\item Studied photodetector/APD noise, laser thermal control, LiDAR detection, camera sensors and spectrometer calibration through parameter scans and uncertainty analysis.
\end{highlights}
\end{onecolentry}
\vspace{0.08 cm}
\section{Work Experience and Career Break}
\begin{twocolentry}{\textit{Dec 2024 - Feb 2025}}
{\large\textbf{Python Developer | BaryonBib}}

\textit{MAX!mize start-up team, MPA, Garching}
\end{twocolentry}
\begin{onecolentry}
\begin{highlights}
\item Contributed to a centralised bibliography platform in Python; implemented and tested LLM-based functionality.
\end{highlights}
\end{onecolentry}
\begin{twocolentry}{\textit{Jan 2025 - Dec 2025}}
\textbf{Medical leave}
\end{twocolentry}
\begin{onecolentry}
Recovery and rehabilitation following an accident and surgery.
\end{onecolentry}
\section{Talks, Workshops and Schools}
\begin{twocolentry}{\textit{Sep 2026}}
\textbf{MAVIS 2026 workshop | Participant}
\end{twocolentry}
\vspace{0.06 cm}
\begin{twocolentry}{\textit{Dec 2024}}
\textbf{Tonale Winter School on Cosmology | Participant}
\end{twocolentry}
\vspace{0.06 cm}
\begin{twocolentry}{\textit{Nov 2024}}
\textbf{ORIGINS Lensing Day}
\end{twocolentry}
\begin{onecolentry}
\href{https://indico.ph.tum.de/event/7717/timetable/?view=standard\#12-watershed-on-the-curved-sky}{Contributed talk: Watershed on the curved sky}
\end{onecolentry}
\vspace{0.06 cm}
\begin{twocolentry}{\textit{Dec 2023}}
\textbf{Cambridge-LMU Cosmology Meeting | Participant}
\end{twocolentry}
\vspace{0.06 cm}
\begin{twocolentry}{\textit{Nov 2022}}
\textbf{Advanced Astrophysical Seminar | LMU Munich}
\end{twocolentry}
\begin{onecolentry}
40-minute course presentation and literature study on large ground-based astronomical telescopes, telescope designs and adaptive optics.
\end{onecolentry}
\vspace{0.06 cm}
\section{Technical Skills}
\begin{onecolentry}
\textbf{Optical modelling:} Wavefront sensing, Fourier PSFs, calibration, regularised reconstruction and closed-loop diagnostics; HCIPy and TIPTOP; Zemax through optical-design coursework.
\end{onecolentry}
\vspace{0.04 cm}
\begin{onecolentry}
\textbf{Scientific computing:} Python (NumPy, SciPy, Matplotlib, pandas, scikit-learn); MATLAB; Mathematica; statistical analysis and MCMC.
\end{onecolentry}
\vspace{0.04 cm}
\begin{onecolentry}
\textbf{Experimental methods:} CAD, 3D printing, mechanical assembly, electronic measurement integration, optical alignment and interferometry.
\end{onecolentry}
\vspace{0.04 cm}
\begin{onecolentry}
\textbf{Research tools:} Linux, Git, Jupyter and LaTeX.
\end{onecolentry}
\vspace{0.04 cm}
\section{Languages and Interests}
\begin{onecolentry}
\textbf{Languages:} Chinese (native); English (fluent); German (fluent); French (beginner).
\end{onecolentry}
\begin{onecolentry}
\textbf{Interests:} Astronomical photography and alpine mountaineering.
\end{onecolentry}
\end{document}
